\documentclass[10pt]{article}
\usepackage[utf8]{inputenc}
\usepackage[T1]{fontenc}
\usepackage{graphicx}
\usepackage[export]{adjustbox}
\graphicspath{ {./images/} }
\usepackage{amsmath}
\usepackage{amsfonts}
\usepackage{amssymb}
\usepackage{mhchem}
\usepackage{stmaryrd}
\usepackage{bbold}

\title{Jinchao Xu }


\author{Deep Learning Algorithms and\\
Analysis}
\date{}


\begin{document}
\maketitle

Summer 2020

\includegraphics[max width=\textwidth]{2022_01_05_5a7d88a0a7ec0275e0bbg-02}

\section{Contents}
1 Monte Carlo Methods $\ldots \ldots \ldots \ldots \ldots \ldots \ldots \ldots . . . . . . . . . . . . . . . . . . . . . . . .$

1.1 Monte Carlo methods $\ldots \ldots \ldots \ldots \ldots \ldots \ldots \ldots \ldots$

$1.1 .1 \quad$ A basic result $\ldots \ldots . . . . . . . . . . . . . . . . . . . . . . . . . . .$

$1.1 .2 \quad$ Application $\ldots \ldots \ldots \ldots . . . . . . . . . . . . . . . . . . . .$

$1.2$ Integral representations of functions $\ldots \ldots \ldots \ldots \ldots \ldots \ldots \ldots$

1.2.1 Fourier representation $\ldots \ldots \ldots \ldots \ldots \ldots \ldots \ldots$

1.2.2 Double Fourier representation $\ldots \ldots \ldots \ldots \ldots \ldots \ldots \ldots$

References $\ldots \ldots \ldots \ldots \ldots \ldots \ldots \ldots \ldots \ldots \ldots \ldots \ldots$

\includegraphics[max width=\textwidth]{2022_01_05_5a7d88a0a7ec0275e0bbg-04}

\section{Monte Carlo Methods}
\subsection{Monte Carlo methods}
Let $\lambda \geq 0$ be a probability density function on $G \subset \mathbb{R}^{d}$ such that
$$
\int_{G} \lambda(\omega) d \omega=1 .
$$
For example:
$$
\lambda(\omega)=\frac{1}{|G|}
$$
if $G$ is bounded. The expectation is defined:
$$
\mathbb{E} g:=\int_{G} g(\omega) \lambda(\omega) d \omega
$$
and for any $h=h\left(\omega_{1}, \omega_{2}, \ldots, \omega_{n}\right): G \times G \cdots G \mapsto \mathbb{R}$
$$
\mathbb{E}_{n} h:=\int_{G \times G \times \ldots \times G} h\left(\omega_{1}, \omega_{2}, \ldots, \omega_{n}\right) \lambda\left(\omega_{1}\right) \lambda\left(\omega_{2}\right) \ldots \lambda\left(\omega_{n}\right) d \omega_{1} d \omega_{2} \ldots d \omega_{n}
$$

\subsubsection{A basic result}
Lemma 1. For any $g \in L^{\infty}(G)$, we have
$$
\mathbb{E}_{n}\left(\mathbb{E} g-\frac{1}{n} \sum_{i=1}^{n} g\left(\omega_{i}\right)\right)^{2}=\left\{\begin{array}{c}
\frac{1}{n} \mathbb{E}\left((\mathbb{E} g-g)^{2}\right) \leq \frac{1}{n} \sup _{\omega, \omega \in G}\left|g(\omega)-g\left(\omega^{\prime}\right)\right|^{2} \\
\frac{1}{n}\left(\mathbb{E}\left(g^{2}\right)-(\mathbb{E}(g))^{2}\right) \leq \frac{1}{n} \mathbb{E}\left(g^{2}\right) \leq \frac{1}{n}\|g\|_{L^{\infty}}^{2},
\end{array}\right.
$$
Proof. First note that
$$
\begin{aligned}
\left(\mathbb{E} g-\frac{1}{n} \sum_{i=1}^{n} g\left(\omega_{i}\right)\right)^{2} &=\frac{1}{n^{2}}\left(n \mathbb{E} g-\sum_{i=1}^{n} g\left(\omega_{i}\right)\right)^{2}=\frac{1}{n^{2}}\left(\sum_{i=1}^{n}\left(\mathbb{E} g-g\left(\omega_{i}\right)\right)\right)^{2} \\
&=\frac{1}{n^{2}} \sum_{i, j=1}^{n}\left(\mathbb{E} g-g\left(\omega_{i}\right)\right)\left(\mathbb{E} g-g\left(\omega_{j}\right)\right) \\
&=\frac{I_{1}}{n^{2}}+\frac{I_{2}}{n^{2}}
\end{aligned}
$$
with

(1.7) $\left.I_{1}=\sum_{i=1}^{n}\left(\mathbb{E} g-g\left(\omega_{i}\right)\right)^{2}, \quad I_{2}=\sum_{i \neq j}^{n}\left((\mathbb{E} g)^{2}-\mathbb{E}(g)\left(g\left(\omega_{i}\right)+g\left(\omega_{j}\right)\right)+g\left(\omega_{i}\right) g\left(\omega_{j}\right)\right)\right)$

Consider $I_{1}$, for any $i$,
$$
\mathbb{E}_{n}\left(\mathbb{E} g-g\left(\omega_{i}\right)\right)^{2}=\mathbb{E}_{n}(\mathbb{E} g-g)^{2}
$$
Thus,
$$
\mathbb{E}_{n} I_{1}=n \mathbb{E}\left((\mathbb{E} g-g)^{2}\right)
$$
For $I_{2}$, note that
$$
\mathbb{E}_{n} g\left(\omega_{i}\right)=\mathbb{E}_{n} g\left(\omega_{j}\right)=\mathbb{E}(g)
$$
and, for $i \neq j$,
$$
\begin{aligned}
\mathbb{E}_{n}\left(g\left(\omega_{i}\right) g\left(\omega_{j}\right)\right) &=\int_{G \times G \times \ldots \times G} g\left(\omega_{j}\right) g\left(\omega_{j}\right) \lambda\left(\omega_{1}\right) \lambda\left(\omega_{2}\right) \ldots \lambda\left(\omega_{n}\right) d \omega_{1} d \omega_{2} \ldots d \omega_{n} \\
&=\int_{G \times G} g\left(\omega_{j}\right) g\left(\omega_{j}\right) \lambda\left(\omega_{1}\right) \lambda\left(\omega_{1}\right) \lambda\left(\omega_{2}\right) d \omega_{1} d \omega_{2} \\
&=\mathbb{E}_{n}\left(g\left(\omega_{i}\right)\right) \mathbb{E}_{n}\left(g\left(\omega_{j}\right)\right)=[\mathbb{E}(g)]^{2}
\end{aligned}
$$
Thus
$$
\mathbb{E}_{n}\left(I_{2}\right)=\mathbb{E}_{n}\left(\sum_{i \neq j}^{n}\left((\mathbb{E} g)^{2}-\mathbb{E}(g)\left(\mathbb{E}\left(g\left(\omega_{i}\right)\right)+\mathbb{E}\left(g\left(\omega_{j}\right)\right)\right)+\mathbb{E}\left(g\left(\omega_{i}\right) g\left(\omega_{j}\right)\right)\right)\right)=0
$$
Consequently, there exist the following two formulas for $\mathbb{E}_{n}\left(\mathbb{E} g-\frac{1}{n} \sum_{i=1}^{n} g\left(\omega_{i}\right)\right)^{2}:$
$$
\mathbb{E}_{n}\left(\mathbb{E} g-\frac{1}{n} \sum_{i=1}^{n} g\left(\omega_{i}\right)\right)^{2}=\frac{1}{n^{2}} \mathbb{E}_{n} I_{1}=\left\{\begin{array}{c}
\frac{1}{n} \mathbb{E}\left((\mathbb{E} g-g)^{2}\right) \\
\frac{1}{n}\left(\mathbb{E}\left(g^{2}\right)-(\mathbb{E} g)^{2}\right)
\end{array}\right.
$$
Based on the first formula above, since
$$
|g(\omega)-\mathbb{E} g|=\left|\int_{G}(g(\omega)-g(\tilde{\omega})) \lambda(\tilde{\omega}) d \tilde{\omega}\right| \leq \sup _{\omega, \omega \in G}\left|g(\omega)-g\left(\omega^{\prime}\right)\right|
$$
it holds that
$$
\mathbb{E}_{n}\left(\mathbb{E} g-\frac{1}{n} \sum_{i=1}^{n} g\left(\omega_{i}\right)\right)^{2} \leq \frac{1}{n} \sup _{\omega, \omega \in G}\left|g(\omega)-g\left(\omega^{\prime}\right)\right|^{2}
$$
Due to the second formula above,
$$
\mathbb{E}_{n}\left(\mathbb{E} g-\frac{1}{n} \sum_{i=1}^{n} g\left(\omega_{i}\right)\right)^{2} \leq \frac{1}{n} \mathbb{E}\left(g^{2}\right) \leq \frac{1}{n}\|g\|_{L^{\infty}}^{2}
$$
which completes the proof. $\square$

Of course, we can use this to prove a high probability result.

Corollary 1. Under the assumptions of the preceding lemma, we have
$$
\overline{\mathbb{P}}\left[\left(\mathbb{E} g-\frac{1}{n} \sum_{i=1}^{n} g\left(\omega_{i}\right)\right)^{2}>\frac{k}{n}\|g\|_{L^{\infty}}^{2}\right]<\frac{1}{k}
$$
Proof.
$$
\overline{\mathbb{P}}\left[\left(\mathbb{E} g-\frac{1}{n} \sum_{i=1}^{n} g\left(\omega_{i}\right)\right)^{2}>\epsilon\right] \leq \epsilon^{-1} \overline{\mathbb{E}}\left(\mathbb{E} g-\frac{1}{n} \sum_{i=1}^{n} g\left(\omega_{i}\right)\right)^{2} \leq \frac{1}{n \epsilon}\|g\|_{L^{\infty}}^{2}
$$
$\square$

This corollary implies that the set of $\omega_{i}$ where the estimate $n^{-1} \sum_{i=1}^{n} g\left(\omega_{i}\right)$ is far from the desired value $\mathbb{E} g$ is small.

The practical usefulness of this algorithm depends upon the existence of a repeatable process (for instance some physical process) which generates $\omega$ according to a desired distribution $\mu$.

The precise meaning of this last statement is essentially that the strong law of large numbers holds. Specifically, if $\omega_{1}, \ldots, \omega_{n}, \ldots$ is a infinite sequence generated by the process, and $A \subset \Omega$ is any a measurable set, then
$$
\lim _{n \rightarrow \infty} \frac{1}{n} \sum_{i=1}^{n} \chi_{A}\left(\omega_{i}\right)=\mu(A) .
$$
Generating $n$ independent samples means generating $\omega_{1}, \ldots, \omega_{n}$ from $\mu^{n}$ according to the above notion. The existence of a realizable process generating samples from a probability distribution, and the practical use of such processes is an interesting topic in the intersection of statistics, physics, and computer science. In addition, statistics/probability theory studies how to take samples from one probability distribution and transform them to samples from another distribution.

\subsubsection{Application}
\section{Lemma 2. Let}
$$
f(x)=\int_{G} g(x, \theta) \lambda(\theta) d \theta=\mathbb{E}(g)
$$
with $\lambda(\theta) \geq 0$ and $\|\lambda(\theta)\|_{L^{1}(G)}=1$. For any $n \geq 1$, there exist $\theta_{i}^{*} \in G$ such that
$$
\left\|f-f_{n}\right\|_{L^{2}(\Omega)}^{2} \leq \frac{1}{n} \int_{G}\|g(\cdot, \theta)\|_{L^{2}(\Omega)}^{2} \lambda(\theta) d \theta=\frac{1}{n} \mathbb{E}\left(\|g(\cdot, \theta)\|_{L^{2}(\Omega)}^{2}\right)
$$
where $\|g(\cdot, \theta)\|_{L^{2}(\Omega)}^{2}=\int_{\Omega}[g(x, \theta)]^{2} d \mu(x)$, and
$$
f_{n}(x)=\frac{1}{n} \sum_{i=1}^{n} g\left(x, \theta_{i}^{*}\right)
$$
Proof. Introducing a probability distribution $\lambda(\theta)$,
$$
f(x)=\mathbb{Z}(g)
$$
By Lemma 1,
$$
\left.\mathbb{E}_{n}\left(\left(\mathbb{E}(g(x, \cdot))-\frac{1}{n} \sum_{i=1}^{n} g\left(x, \theta_{i}\right)\right)\right)^{2}\right) \leq \frac{1}{n} \mathbb{E}\left(g^{2}\right)
$$
and
$$
\mathbb{E}_{n}\left(h\left(\theta_{1}, \theta_{2}, \cdots, \theta_{n}\right)\right) \leq \frac{1}{n} \mathbb{E}\left(\int_{\Omega} g^{2} d \mu(x)\right),
$$
by taking integral where
$$
\left.h\left(\theta_{1}, \theta_{2}, \cdots, \theta_{n}\right)=\int_{\Omega}\left(\mathbb{E}(g(x, \cdot))-\frac{1}{n} \sum_{i=1}^{n} g\left(x, \theta_{i}\right)\right)\right)^{2} d \mu(x)
$$
Sine $\mathbb{E}_{n}(1)=1$ and $\mathbb{E}_{n}(h) \leq \frac{1}{n} \mathbb{E}\left(\int_{O} g^{2} d \mu(x)\right)$, there exists $\left(\theta_{1}^{*}, \theta_{2}^{*}, \cdots, \theta_{n}^{*}\right) \in G \times G \times$ $\cdots \times G$ such that
$$
h\left(\theta_{1}^{*}, \theta_{2}^{*}, \cdots, \theta_{n}^{*}\right) \leq \frac{1}{n} \int_{\Omega} \mathbb{E}\left(g^{2}\right) d \mu(x)
$$
Otherwise, $\mathbb{E}_{n}(h)>\frac{1}{n} \mathbb{E}\left(\int_{\Omega} g^{2} d \mu(x)\right)$ if $\left.h\left(\theta_{1}, \theta_{2}, \cdots, \theta_{n}\right)>\frac{1}{n} \int_{\Omega} \mathbb{E}\left(g^{2}\right)\right) d \mu(x) .$ This implies that
$$
\left\|f-f_{n}\right\|_{L^{2}(\Omega)}^{2} \leq \frac{1}{n} \int_{G}\|g(\cdot, \theta)\|_{L^{2}(\Omega)}^{2} \lambda(\theta) d \theta,
$$
which completes the proof.

We also have a more general version of the above lemma.

\section{Lemma 3. Let}
$$
f(x)=\int_{G} g(x, \theta) \lambda(\theta) d \theta=\mathbb{E}(g)
$$
with $\|\lambda(\theta)\|_{L^{1}(\Theta)}=1$. For any $n \geq 1$, there exist $\theta_{i}^{*} \in G$ such that
$$
\left\|f-f_{n}\right\|_{H^{m}(\Omega)}^{2} \leq \int_{G}\|g(\cdot, \theta)\|_{H^{m}(\Omega)}^{2} \lambda(\theta) d \theta=\frac{1}{n} \mathbb{E}\left(\|g(\cdot, \theta)\|_{H^{m}(\Omega)}^{2}\right)
$$
where
$$
f_{n}(x)=\frac{1}{n} \sum_{i=1}^{n} g\left(x, \theta_{i}^{*}\right)
$$
In particular, if
$$
\left|D^{\alpha} g(x, \theta)\right| \leq C, \quad \forall x, \theta,|\alpha| \leq m
$$
Then
$$
\left\|f-f_{n}\right\|_{H^{m}(\Omega)} \leq\left(\begin{array}{c}
m+d \\
m
\end{array}\right)^{1 / 2}|\Omega|^{1 / 2} n^{-1 / 2}
$$
For any $f(x)=\int_{G} g(x, \theta) \rho(\theta) d \theta$ with $\|\rho\|_{L^{1}(\Theta)} \neq 1 .$ Let $\lambda(\theta)=\frac{\rho(\theta)}{\|\rho\|_{L^{1}(\theta)}} .$ Thus,
$$
f(x)=\|\rho\|_{L^{1}(\Theta)} \int_{G} g(x, \theta) \lambda(\theta) d \theta
$$
with $\|\lambda(\theta)\|_{L^{1}(\Theta)}=1 .$ We can apply the above two lemmas to the given function $f(x)$.

\subsection{Integral representations of functions}
\subsubsection{Fourier representation}
Consider the Fourier transform:
$$
\hat{f}(\omega)=\frac{1}{(2 \pi)^{d}} \int_{\mathbb{R}^{d}} e^{-i \omega \cdot x} f(x) d x \quad \forall \omega \in \mathbb{R}^{d}
$$
We write $\hat{f}(\omega)=e^{i \theta(\omega)}|\hat{f}(\omega)| .$ By Fourier inversion formula,
$$
f(x)=\int_{\mathbb{R}^{d}} e^{i \omega \cdot x} \hat{f}(\omega) d \omega=\int_{\mathbb{R}^{d}} e^{i(\omega \cdot x+\beta(\omega))}|\hat{f}(\omega)| d \omega .
$$
Since $f(x)$ is real-valued, it implies that, for $x$
$$
\begin{aligned}
f(x) &=\operatorname{Re} \int_{\mathbb{R}^{d}} e^{i \omega \cdot x} \hat{f}(\omega) d \omega \\
&=\operatorname{Re} \int_{\mathbb{R}^{d}} e^{i \omega \cdot x} e^{i \beta(\omega)}|\hat{f}(\omega)| d \omega \\
&=\int_{\mathbb{R}^{d}} \cos (\omega \cdot x+\beta(\omega))|\hat{f}(\omega)| d \omega
\end{aligned}
$$
Then we have
$$
f(x)=\int_{\mathbb{R}^{d}} k(x, \omega) d \omega
$$
with
$$
k(x, \omega)=\cos (\omega \cdot x+\beta(\omega))|\hat{f}(\omega)|
$$
and
$$
|k(x, \omega)| \leq|\hat{f}(\omega)|=\rho(\omega)
$$
Theorem 1. There exist $\omega_{i} \in \mathbb{R}^{d}$, s.t., $G=\mathbb{R}$ and
$$
\int_{\Omega}\left(f(x)-f_{n}(x)\right)^{2} \leq \frac{1}{n} \int_{\mathbb{R}^{d}}|\hat{f}(\omega)| d \omega,
$$
where
$$
f_{n}(x)=\frac{\|\hat{f}\|_{L^{1}}}{n} \sum_{i=1}^{n} \frac{\cos \left(\omega_{i}^{*} \cdot x+\beta_{i}^{*}\right)}{\rho\left(\omega_{i}^{*}\right)}
$$
Note that
$$
f_{n}=\sum_{i=1}^{n} \frac{\cos \left(\omega_{i}^{*} \cdot x+\beta_{i}^{*}\right)}{\rho\left(\omega_{i}^{*}\right)} \in{ }_{n} \mathrm{~N}(\sigma, n)
$$
with
$$
\sigma(t)=\cos (t)
$$

\subsubsection{Double Fourier representation}
Assume that $\sigma$ is a locally Riemann integrable function and $\sigma \in L^{1}(\mathbb{R})$ and thus the Fourier transform of $\sigma$ is well-defined and continuous. Since $\sigma$ is non-zero and
$$
\hat{\sigma}(\omega)=\frac{1}{2 \pi} \int_{\mathbb{R}} \sigma(t) e^{-i \omega t} d t
$$
this implies that $\hat{\sigma}(a) \neq 0$ for some $a \neq 0$. Via a change of variables $t=w \cdot x+b$ and $d t=d b$, this means that for all $x$ and $\omega$, we have
$$
\begin{aligned}
0 \neq \hat{\sigma}(a) &=\frac{1}{2 \pi} \int_{\mathbb{R}} \sigma(\omega \cdot x+b) e^{-i a(\omega \cdot x+b)} d b \\
&=e^{-i a \omega \cdot x} \frac{1}{2 \pi} \int_{\mathbb{R}} \sigma(\omega \cdot x+b) e^{-i a+b} d b
\end{aligned}
$$
and so
$$
e^{i a \omega \cdot x}=\frac{1}{2 \pi \hat{\sigma}(a)} \int_{\mathbb{R}} \sigma(\omega \cdot x+b) e^{-i a b} d b
$$
Likewise, since the growth condition also implies that $\sigma^{(k)} \in L^{1}$, we can differentiate the above expression under the integral with respect to $x$.

This allows us to write the Fourier mode $e^{i a \omega \cdot x}$ as an integral of neuron output functions. We substitute this into the Fourier representation of $f$ (note that the assumption we make implies that $\hat{f} \in L^{1}$ so this is rigorously justified for a.e. $x$ ) to get

$f(x)=\int_{\mathbb{R}^{d}} e^{i \omega \cdot x} \hat{f}(\omega) d \omega=\int_{\mathbb{R}^{d}} \int_{\mathbb{R}} \frac{1}{2 \pi \hat{\sigma}(a)} \sigma\left(a^{-1} \omega \cdot x+b\right) \hat{f}(\omega) e^{-i a b} d b d \omega=\int_{\mathbb{R}^{d} \times \mathbb{R}} k(x, \theta) d \theta$ where $\theta=(\omega, b)$ and
$$
k(x, \theta)=\frac{1}{2 \pi \hat{\sigma}(a)} \sigma\left(a^{-1} \omega \cdot x+b\right) \hat{f}(\omega) e^{-i a b}
$$
Thus we have
$$
\begin{aligned}
|k(x, \theta)| & \leq \frac{1}{2 \pi|\hat{\sigma}(a)|} \max _{x \in \Omega}\left|\sigma\left(a^{-1} \omega \cdot x+b\right) \| \hat{f}(\omega)\right| \\
& \leq h(\omega, b)|\hat{f}(\omega)|=\rho(\theta)
\end{aligned}
$$
where
$$
h(\omega, b)=\max _{x \in \Omega}\left|\sigma\left(a^{-1} \omega \cdot x+b\right)\right|
$$
if we ignore the coefficient. Thus, following the discussion in last section, the next step is to analyze $h(\omega, b)$ which we will discuss in the next section. Before that, let us introduce a special case of the above representation once the activation function is periodic.

\includegraphics[max width=\textwidth]{2022_01_05_5a7d88a0a7ec0275e0bbg-12}

\section{References}
[1] L. Chen, P. Sun, and J. Xu. Optimal anisotropic meshes for minimizing interpolation errors in $l^{p}$-norm. Mathematics of Computation, $76(257): 179-204,2007 .$

[2] K. He, X. Zhang, S. Ren, and J. Sun. Deep residual learning for image recognition. In Proceedings of the IEEE Conference on Computer Vision and Pattern Recognition, pages 770-778, 2016

[3] K. He, X. Zhang, S. Ren, and J. Sun. Identity mappings in deep residual networks. In European Conference on Computer Vision, pages 630-645. Springer, $2016 .$

[4] M. Leshno, V.Y. Lin, A. Pinkus, and S. Schocken. Multilayer feedforward networks with a nonpolynomial activation function can approximate any function. Neural networks, 6(6):861-867, 1993 .

[5] A. Pinkus. Approximation theory of the mlp model in neural networks. Acta numerica, $8: 143-195,1999 .$

[6] L. R. Scott and S. Zhang. Finite element interpolation of nonsmooth functions satisfying boundary conditions. Mathematics of Computation, $54(190): 483-493$, $1990 .$

[7] J. Xu. Finite element methods. Lecture notes, 2020 .


\end{document}